\subsection{Project Context}

\ChapterContent{
  The integration of digital infrastructure within educational administrative sectors has transitioned from a luxury to a fundamental necessity for operational efficiency. Modern institutions leverage automated systems to handle complex workflows, ranging from student records to logistical resource management (Chen \& Patel, 2024). By reducing reliance on manual data entry, these technologies minimize human error and ensure that administrative staff can focus on high-level strategic planning rather than repetitive clerical tasks. \\

  Global research highlights a significant trend toward the adoption of Cloud-based Enterprise Resource Planning (ERP) systems to manage institutional assets. Studies conducted in European and North American universities demonstrate that real-time tracking systems can reduce asset loss by up to 30\% within the first year of implementation (Schmidt, 2023). Furthermore, the shift toward mobile-integrated solutions, such as IoT and scanning technologies, has become the international standard for maintaining transparent audit trails in large-scale academic environments.\\

  In the local context, many educational institutions are currently undergoing a digital transformation to meet national accreditation standards for resource management. Recent local literature suggests that while many schools have digitized their student portals, the ``back-office'' inventory processes often remain neglected or semi-manual (Reyes \& Santos, 2025). These studies emphasize that localized IT solutions are critical in overcoming budget constraints while improving the accountability of government-funded equipment and facilities.\\

  X University is a leading academic institution dedicated to providing quality education across multiple campuses and specialized laboratory facilities. Currently, the university's inventory process relies on a traditional ledger-based system supplemented by disconnected departmental spreadsheets. When new equipment is acquired, it is manually tagged with a physical sticker and logged into a central file, which is only updated during the annual year-end audit.\\

  To address these inefficiencies, the Inventory System with QR Code Generation is proposed as a comprehensive digital solution for X University. This system will automate the tagging process by generating unique, scan-able QR codes for every asset, allowing for instant data retrieval via mobile devices. By centralizing the database, the output aims to provide real-time visibility into the university's resources, streamlining the audit process and ensuring better asset accountability.\\
}